%! AP Statistics — Unit 2, Lesson 2.3 — Warm-Up (Answer Key)
\documentclass[10pt]{article}
\usepackage{apstats-article}
\usepackage{apstats-key}

%=============================================================================
\begin{document}

\pageheader{Unit 2, Lesson 2.3}{Warm-Up --- Answer Key}
\begin{tabularx}{\linewidth}{X X X}
Name:\;\ans{Key} & Date:\; & Period:\;
\end{tabularx}

\vspace{0.3in}

{\large\bfseries\color{navy} Part 1 \quad Conditional Relative Frequencies (Lesson 2.2 Review)}

\vspace{0.12in}

\noindent A survey of 200 high school students recorded two variables: whether they play
a sport (Yes / No) and whether they have a part-time job (Yes / No). Results are below.

\vspace{0.1in}

\renewcommand{\arraystretch}{1.9}
\begin{tabularx}{\linewidth}{@{} l C{2.8cm} C{2.8cm} C{2.4cm} @{}}
\rowcolor{sky}
                    & \textbf{Job: Yes} & \textbf{Job: No} & \textbf{Row Total} \\
\hline
\textbf{Sport: Yes} & 30 & 90 & \ans{120} \\
\textbf{Sport: No}  & 40 & 40 & \ans{80} \\
\hline
\textbf{Col.\ Total}& \ans{70} & \ans{130} & \ans{200} \\
\end{tabularx}

\vspace{0.18in}

\begin{enumerate}

    \item \ans{Row totals: 120, 80. Column totals: 70, 130. Grand total: 200.}

    \vspace{0.06in}

    \item $P(\text{Job: Yes} \mid \text{Sport: Yes}) = \dfrac{\ans{30}}{\ans{120}} = \ans{0.25 = 25\%}$

    \vspace{0.18in}

    \item $P(\text{Job: Yes} \mid \text{Sport: No}) = \dfrac{\ans{40}}{\ans{80}} = \ans{0.50 = 50\%}$

    \vspace{0.14in}

    \item Circle: \textbf{Yes} \textcolor{keyred}{\textbf{$\leftarrow$ correct}}\\[8pt]
    \ans{Student-athletes have a job rate of 25\%, while non-athletes have a job rate of 50\% ---
    a substantial difference, suggesting an association between playing a sport and having a job.}

\end{enumerate}

\vspace{0.28in}

{\large\bfseries\color{navy} Part 2 \quad Putting a Number on the Difference}

\vspace{0.12in}

\begin{tcolorbox}[colback=greenbg, colframe=greenacc, boxrule=0.8pt, arc=2mm,
    left=4mm, right=4mm, top=3mm, bottom=3mm]
    \small
    \textbf{Think about it:} In Part 1 you found two conditional proportions and compared
    them informally. Today we will compute \emph{statistics} that quantify exactly
    \emph{how different} the two proportions are.
\end{tcolorbox}

\vspace{0.14in}

\noindent $\hat{p}_{\text{athlete}} - \hat{p}_{\text{non-athlete}} = \ans{0.25} - \ans{0.50} = \ans{-0.25}$

\vspace{0.1in}

\noindent \ansline{Student-athletes are 25 percentage points \emph{less} likely to have a job
than non-athletes.}

\begin{teachernote}
\textbf{Item 4:} Many students will circle ``No'' because the raw job counts for non-athletes
(40 vs.\ 40) look balanced. Emphasize conditioning: the relevant comparison is 25\% vs.\ 50\%,
not 30 vs.\ 40.

\textbf{Part 2 sign:} The sign is negative because athletes are the first group. Students who
flip the subtraction get $+0.25$ --- accept either as long as direction is stated in context.
\end{teachernote}

\end{document}
